\begin{problem}{자리 바꾸기}{standard input}{standard output}{1 second}{512 megabytes}

$N$개의 의자가 원형으로 배치되어 있고, 각 의자에는 한 명씩 총 $N$명의 사람이 앉아 있다. 처음 상태에서 임의의 의자를 $1$번 의자로 정하고, 시계 방향으로 의자에 $1$번부터 $N$번까지 번호를 붙인다. 이때 $i$번 의자에 앉아 있는 사람을 $i$번 사람이라고 하자.

사람들은 자리를 바꾸기 위해 다음과 같은 과정을 반복한다. 모든 사람은 동시에 현재 위치에서 시계 방향으로 자신의 번호만큼 떨어진 의자로 이동한다. 이동이 끝난 뒤 모든 의자에 정확히 한 명씩 위치해 있다면, 그 상태로 자리를 확정하여 앉는다. 그렇지 않다면 아무도 자리를 확정하지 않고, 이동이 끝난 위치에서 같은 과정을 다시 반복한다.

이때 위 과정을 몇 번 수행해야 모든 사람이 자리에 앉게 되는지 구하여라.

\InputFile
첫 번째 줄에 의자의 개수 $N(3 \le N \le 10^{18})$이 주어진다.

\OutputFile
첫 번째 줄에 이동이 끝난 뒤 모든 의자에 정확히 한 명씩 위치하게 되는 상태가 처음으로 나타날 때까지 수행해야 하는 과정의 횟수를 출력한다.

\Examples

\begin{example}
\exmpfile{example.01}{example.01.a}%
\exmpfile{example.02}{example.02.a}%
\end{example}

\Note
첫 번째 예제의 경우 첫 번째 과정에서 1번은 2번 자리로, 2번은 1번 자리로, 3번은 3번 자리로 이동하여 앉게 된다.

계산 시 32비트 정수 범위를 초과하는 경우가 발생할 수 있음에 유의하여라.

\end{problem}